\documentclass[../../main.tex]{subfiles}
\graphicspath{{../../image/}{../../figure/}}
\begin{document}

\section{解的正则性}\label{section:ch12-解的正则性}

迄今为止收集的信息为解的正则性描绘出如下图景：

\begin{theorem}\label{theorem:ch12-正则性}
设 $U(x,t)$ 是真非线性两守恒律组\eqref{align:ch12-系统}的容许 $BV$ 解，具有前几节所述性质。则
\begin{enumerate}
\item[\textup{(a)}] 近似连续的任一点 $(\bar x,\bar t)$ 是 $U$ 的连续点。
\item[\textup{(b)}] 近似跳跃间断的任一点 $(\bar x,\bar t)$ 是 $U$ 的（经典）跳跃间断点。
\item[\textup{(c)}] 任一不规则点 $(\bar x,\bar t)$ 是某族或两族特征的中心压缩波的聚焦点，和/或同族或异族激波的相互作用点。
\item[\textup{(d)}] 不规则点集（至多）可数。
\end{enumerate}
\end{theorem}

\begin{proof}
断言 (a)、(b)、(c) 是\cref{theorem:ch12-局部正则存在}的推论。特别地，$(\bar x,\bar t)$ 是近似连续点当且仅当 $(z_W,w_W)=(z_E,w_E)$，此时四个极限 $(z_W,w_W)$、$(z_E,w_E)$、$(z_N,w_N)$、$(z_S,w_S)$ 全部重合。当 $(z_W,w_W)\neq(z_E,w_E)$ 时，$(\bar x,\bar t)$ 是 1-激波集的近似跳跃间断点当 $(z_W,w_W)=(z_S,w_S)$、$(z_E,w_E)=(z_N,w_N)$；或是 2-激波集的近似跳跃间断点当 $(z_W,w_W)=(z_N,w_N)$、$(z_E,w_E)=(z_S,w_S)$；而在所有其他情形为不规则点。

为验证断言 (d)，设不规则点 $I=(\bar x,\bar t)$ 是某 1-特征线树 $\mathscr{M}$ 或 2-特征线树 $\mathscr{N}$ 的节点。若 $I$ 是中心 1-压缩波的聚焦点和/或 1-激波的相互作用点，则由\eqref{align:ch12-P1}，$I$ 将对 $\mathscr{P}_1(\mathscr{M})$ 贡献正值。类似地，若 $I$ 是中心 2-压缩波的聚焦点和/或 2-激波的相互作用点，则由\eqref{align:ch12-P2}，$I$ 将对 $\mathscr{P}_2(\mathscr{N})$ 贡献正值。最后，设 $I$ 是 1-激波与 2-激波的相互作用点。我们在 $\chi_-^I$ 图像上靠近 $I$ 处向 $\mathscr{M}$ 附加一额外节点 $K$。于是 $|\Delta w^{KI}-\Delta w_*^I|>0$，从而由\eqref{align:ch12-Q1}，我们得到对 $\mathscr{Q}_1(\mathscr{M})$ 的正贡献。由于波相互作用总量有界，由\cref{lemma:ch12-相互作用有界}，我们断定不规则点集必然（至多）可数。证明完成。
\end{proof}

定理 11.3.5 的类比在此同样成立：

\begin{theorem}\label{theorem:ch12-局部Lipschitz}
设解 $U$ 的连续点集 $\mathscr{C}$ 有非空内部 $\mathscr{C}^0$。则 $U$ 在 $\mathscr{C}^0$ 上局部 Lipschitz 连续。
\end{theorem}

\begin{proof}
我们验证 $z$ 在 $\mathscr{C}^0$ 上局部 Lipschitz 连续。设 $(\bar x,\bar t)\in\mathscr{C}^0$，且 $\mathscr{C}$ 含矩形 $\{(x,t):|x-\bar x|<kp,\ |t-\bar t|<p\}$，其中 $p>0$，$k$ 相对于 $|\lambda|$ 与 $\mu$ 充分大。平移坐标轴，不妨设 $\bar t=p$。我们固定 $\bar y>\bar x$，使 $\bar y-\bar x$ 相对于 $p$ 充分小，并以 $x_\ell=\bar x$、$x_r=\bar y$ 应用\eqref{align:ch12-散开估计1}。由于解在矩形内连续，$\mathscr{P}_2(\mathscr{F})$ 与 $\mathscr{Q}_2(\mathscr{F})$ 均为零，故回忆\eqref{align:ch12-barV}，
\begin{align}
z(\bar x,\bar t)-z(\bar y,\bar t)\leqslant\frac{2\hat c}{p}(\bar y-\bar x).
\label{align:ch12-Lipschitz1}
\end{align}

函数 $(\hat z,\hat w)(x,t)=(z,w)(x+\bar y-x,2p-t)$ 是另一解 $\hat U$ 的 Riemann 不变量，该解在矩形 $\{(x,t):|x-\bar y|<kp,\ |t-\bar t|<p\}$ 上连续从而容许。将\eqref{align:ch12-Lipschitz1}应用于 $\hat z$，得
\begin{align}
z(\bar y,\bar t)-z(\bar x,\bar t)=\hat z(\bar x,\bar t)-\hat z(\bar y,\bar t)\leqslant\frac{2\hat c}{p}(\bar y-\bar x).
\label{align:ch12-Lipschitz2}
\end{align}

现固定 $\bar s>\bar t$，使 $\bar s-\bar t$ 相对于 $p$ 充分小。我们构造从 $(\bar x,\bar s)$ 发出的极小后向 1-特征线 $\xi$，它在点 $\bar y=\xi(\bar t)$ 处被 $\bar t$-时间线截取，其中 $0<\bar y-\bar x\leqslant-\bar\lambda(\bar s-\bar t)$。由\cref{theorem:ch12-传播主定理}，$z(\bar x,\bar s)=z(\bar y,\bar t)$，故由\eqref{align:ch12-Lipschitz1}与\eqref{align:ch12-Lipschitz2}，
\begin{align}
|z(\bar x,\bar s)-z(\bar x,\bar t)|\leqslant\frac{2\hat c}{p}(\bar y-\bar x)\leqslant-\frac{2\bar\lambda\hat c}{p}(\bar s-\bar t).
\label{align:ch12-Lipschitz3}
\end{align}
于是 $z$ 是 Lipschitz 的。类似论证表明 $w$ 在 $\mathscr{C}^0$ 上也是 Lipschitz 的。证明完成。
\end{proof}

\end{document}
